%% 王彬 — 个人简历（中文版）
%% 编译命令：xelatex cv_zh.tex

\documentclass[10pt, a4paper]{article}

% ─── 宏包 ─────────────────────────────────────────────────────────────────────
\usepackage[a4paper, top=1.8cm, bottom=1.8cm, left=2.0cm, right=2.0cm]{geometry}
\usepackage{hyperref}
\usepackage{fontawesome5}
\usepackage{enumitem}
\usepackage{titlesec}
\usepackage{xcolor}
\usepackage{microtype}
\usepackage{parskip}
\usepackage{ctex}

% ─── 字体设置 ─────────────────────────────────────────────────────────────────
\setCJKmainfont{Noto Serif CJK SC}[
  BoldFont   = Noto Serif CJK SC Bold,
  ItalicFont = Noto Serif CJK SC
]
\setCJKsansfont{Noto Sans CJK SC}
\setCJKmonofont{Noto Sans CJK SC}

% ─── 颜色 ─────────────────────────────────────────────────────────────────────
\definecolor{accent}{RGB}{99, 102, 241}

% ─── 超链接 ───────────────────────────────────────────────────────────────────
\hypersetup{
    colorlinks = true,
    urlcolor   = accent,
    linkcolor  = accent,
    pdftitle   = {王彬 -- 个人简历},
    pdfauthor  = {王彬}
}

% ─── 章节样式 ─────────────────────────────────────────────────────────────────
\titleformat{\section}
    {\normalsize\bfseries}
    {}
    {0em}
    {}
    [\vspace{-0.45em}\rule{\linewidth}{0.6pt}\vspace{0.05em}]
\titlespacing*{\section}{0pt}{0.9em}{0.4em}

% ─── 版式 ─────────────────────────────────────────────────────────────────────
\setlength{\parindent}{0pt}
\setlength{\parskip}{0pt}

% 条目命令：
%   第一行：加粗职位（左）| 时间（右）
%   第二行：斜体单位（左）| 斜体地点（右）
% 用法：\cventry{职位}{单位}{地点}{时间}
\newcommand{\cventry}[4]{%
  \begin{tabular*}{\linewidth}{@{}l@{\extracolsep{\fill}}r@{}}
    \textbf{#1} & {#4} \\[-0.1em]
    \textit{\small #2} & \textit{\small #3}
  \end{tabular*}%
  \vspace{0.1em}%
}

% ─────────────────────────────────────────────────────────────────────────────
\begin{document}
\pagestyle{empty}

% ── 页眉 ──────────────────────────────────────────────────────────────────────
\begin{center}
  {\LARGE\bfseries 王\ 彬}\\[0.35em]
  {\small 华为2012实验室 研究员}\\[0.45em]
  {\small
    \faMapMarker*\ 中国·北京\quad
    \href{https://www.linkedin.com/in/bin-wang-85a164147/}{\faLinkedin\ Bin Wang}\quad
    \href{https://scholar.google.com/citations?user=KWZG_YsAAAAJ&hl=en}{\faGraduationCap\ Google Scholar}\quad
    \href{mailto:binderwang@163.com}{\faEnvelope\ binderwang@163.com}
  }
\end{center}
\vspace{0.3em}

% ── 研究方向 ──────────────────────────────────────────────────────────────────
\section{研究方向}
{\small
我的研究聚焦于构建能够理解、推理并与世界交互的智能系统。
具体方向包括：\textbf{基础模型}（大语言模型、多模态模型、数据合成、后训练与推理），
\textbf{智能体}（由大语言模型驱动的自主智能体，具备规划、工具调用与复杂任务执行能力），
\textbf{强化学习}（策略优化、奖励建模、RLVR \& RLHF 及离线强化学习），
以及\textbf{具身智能}（机器人操作、视觉-语言接地与感觉运动学习）。
}

% ── 工作经历 ──────────────────────────────────────────────────────────────────
\section{工作经历}

\cventry{研究员}{华为2012实验室}{中国·北京}{2018 -- 至今}
\begin{itemize}[leftmargin=1.5em, itemsep=0.1em, topsep=0.15em, parsep=0pt, label=\textbullet]
  \item {\small \textbf{基础模型}（2023--至今）：主导大语言模型/视觉语言模型后训练，聚焦高阶推理能力（ARC）。
    通过合成数据与双系统方法提升复杂推理能力
    （\textit{LogicTree、盘古 Embedded、盘古 Ultra}），以及大语言模型函数调用（\textit{ToolACE}）。}
  \item {\small \textbf{具身智能}（2022--2023）：面向具身智能体的视觉-语言预训练（\textit{EmbodiedGPT}），
    基于大语言模型的任务规划（\textit{Tree-Planner}），扩散模型驱动的机器人操作，
    以及离线元强化学习规划（\textit{MetaDiffuser}）。}
  \item {\small \textbf{自动驾驶}（2018--2022）：强化学习驱动的驾驶决策规划（变道、路口通行）。
    模仿学习（Tripple-GAIL），基于分解互信息的元强化学习（\textit{DOMINO}），
    以及打破 Atari 世界纪录的样本高效行为选择方法（\textit{Learnable Behavior Control}）。}
\end{itemize}

\vspace{0.35em}
\cventry{讲师}{济南大学}{中国·山东}{2015 -- 2018}
\begin{itemize}[leftmargin=1.5em, itemsep=0.1em, topsep=0.15em, parsep=0pt, label=\textbullet]
  \item {\small 从事机器人导航的强化学习研究。}
\end{itemize}

% ── 教育背景 ──────────────────────────────────────────────────────────────────
\section{教育背景}

\cventry{控制科学与工程 博士}%
        {中国科学院自动化研究所}%
        {中国·北京}{2010 -- 2015}
\begin{itemize}[leftmargin=1.5em, itemsep=0.1em, topsep=0.15em, parsep=0pt, label=\textbullet]
  \item {\small 研究方向：强化学习与自动驾驶。}
\end{itemize}

\vspace{0.35em}
\cventry{自动化 本科}{中国石油大学}{中国·青岛}{2006 -- 2010}

% ── 代表性论文 ────────────────────────────────────────────────────────────────
\section{代表性论文}

\begin{enumerate}[leftmargin=*, label={[\arabic*]}, itemsep=0.25em,
                  topsep=0.1em, parsep=0pt]

  \item {\small Zehao Wang, Lin Yang, Jie Wang, Kehan Wang, Hanzhu Chen,
    \textbf{Bin Wang}, Jianye Hao, Defu Lian, Bin Li, Enhong Chen.
    \textit{LogicTree: Improving Complex Reasoning of LLMs via Instantiated
    Multi-step Synthetic Logical Data}. NeurIPS 2025 \textbf{(Spotlight)}.}

  \item {\small Weiwen Liu, Xu Huang, Xingshan Zeng, \ldots,
    \textbf{Bin Wang}, \ldots, Enhong Chen, et al.
    \textit{ToolACE: Winning the Points of LLM Function Calling}. ICLR 2025.}

  \item {\small Yichun Yin, Wenyong Huang, Kaikai Song, \ldots,
    \textbf{Bin Wang}, \ldots, Zhicheng Liu, et al.
    \textit{Pangu Ultra: Pushing the Limits of Dense Large Language Models
    on Ascend NPUs}. arXiv 2025.}

  \item {\small Hanting Chen, Yasheng Wang, Kai Han, \ldots,
    \textbf{Bin Wang}, \ldots, Yunhe Wang, et al.
    \textit{Pangu Embedded: An Efficient Dual-system LLM Reasoner
    with Metacognition}. arXiv 2025.}

  \item {\small Mengkang Hu, Yao Mu, Xinmiao Yu, Mingyu Ding, Shiguang Wu,
    Wenqi Shao, Qiguang Chen, \textbf{Bin Wang}, Yu Qiao, Ping Luo.
    \textit{Tree-Planner: Efficient Close-loop Task Planning with Large
    Language Models}. ICLR 2024.}

  \item {\small Fei Ni, Jianye Hao, Shiguang Wu, Longxin Kou, Jiashun Liu,
    Yan Zheng, \textbf{Bin Wang}, Yuzheng Zhuang.
    \textit{Generate Subgoal Images before Act: Unlocking the Chain-of-Thought
    Reasoning in Diffusion Model for Robot Manipulation with Multimodal Prompts}.
    CVPR 2024.}

  \item {\small Yao Mu, Qinglong Zhang, Mengkang Hu, Wenhai Wang, Mingyu Ding,
    Jun Jin, \textbf{Bin Wang}, Jifeng Dai, Yu Qiao, Ping Luo.
    \textit{EmbodiedGPT: Vision-Language Pre-Training via Embodied Chain
    of Thought}. NeurIPS 2023 \textbf{(Spotlight)}.}

  \item {\small Yao Lai, Jinxin Liu, Zhentao Tang, \textbf{Bin Wang},
    Jianye Hao, Ping Luo.
    \textit{ChiPFormer: Transferable Chip Placement via Offline Decision
    Transformer}. ICML 2023.}

  \item {\small Fei Ni, Jianye Hao, Yao Mu, Yifu Yuan, Yan Zheng,
    \textbf{Bin Wang}, Zhixuan Liang.
    \textit{MetaDiffuser: Diffusion Model as Conditional Planner for
    Offline Meta-RL}. ICML 2023.}

  \item {\small Jiajun Fan, Yuzheng Zhuang, Yuecheng Liu, Jianye Hao,
    \textbf{Bin Wang}, Jiangcheng Zhu, Hao Wang, Shu-Tao Xia.
    \textit{Learnable Behavior Control: Breaking Atari Human World Records
    via Sample-Efficient Behavior Selection}. ICLR 2023 \textbf{(Oral)}.}

  \item {\small Yao Mu, Fei Ni, Weichen Yu, \textbf{Bin Wang},
    Jianye Hao, Ping Luo.
    \textit{DOMINO: Decomposed Mutual Information Optimization for Generalized
    Context in Meta-Reinforcement Learning}.
    NeurIPS 2022 \textbf{(Spotlight)}.}

\end{enumerate}

% ── 荣誉与奖励 ────────────────────────────────────────────────────────────────
\section{荣誉与奖励}

{\small
\begin{tabular*}{\linewidth}{@{}l@{\extracolsep{\fill}}r@{}}
  创新先锋 \& 2012实验室所长奖 & 2025 \\[0.15em]
  2012实验室所长奖             & 2024 \\[0.15em]
  2012实验室所长奖             & 2023 \\[0.15em]
  中央研究院所长团队奖 \& 优秀导师 & 2021 \\[0.15em]
  创新先锋                     & 2020 \\
\end{tabular*}
}

% ── 学术服务 ──────────────────────────────────────────────────────────────────
\section{学术服务}

{\small\textbf{会议审稿人}：NeurIPS、ICLR、ICML 等。}

\end{document}
